
\documentclass[../differential_topology.tex]{subfiles}
\begin{document}
\section{Aula 16 - Índice de Campos Vetoriais e o Teorema de Poincaré--Hopf}
\subsection{Motivações}
\begin{itemize}
	\item Linearização de campos vetoriais;
	\item Teorema de Poincaré--Hopf.
\end{itemize}
\subsection{O Teorema de Poincaré--Hopf}

\begin{theorem*}[Teorema de Poincaré--Hopf]
	Seja M uma variedade compacta e \(v\) um campo vetorial suave em M. Suponhamos que, se \(\partial M\neq\emptyset\), o campo \(v\) aponta para fora de M ao longo de seu bordo. Se suas singularidades são isoladas e
	\[
		\operatorname{Sing}(v)=\{p_1,\dotsc,p_{\ell}\},
	\]
	então
	\[
		\boxed{\chi(M)=\sum_{j=1}^{\ell}i(v,p_j),}
	\]
	onde \(\chi(M)\) é a característica de Euler de M.
\end{theorem*}
\begin{proof*}
	\textbf{Caso inicial:} suponhamos que \(M\subseteq\mathbb{R}^{m}\) seja uma variedade compacta de dimensão m, com bordo. Seja
	\[
		g:\partial M\rightarrow\mathbb{S}^{m-1}
	\]
	a aplicação normal de Gauss, que associa a cada ponto do bordo seu vetor normal exterior unitário. A primeira etapa da demonstração consiste em provar que
	\[
		\boxed{\sum_{j=1}^{\ell}i(v,p_j)=\mathrm{deg}(g).}
	\]

	Escolha \(\varepsilon>0\) suficientemente pequeno para que as bolas fechadas
	\[
		B_i=\overline{B_{\varepsilon}(p_i)}
	\]
	estejam contidas no interior de M e sejam duas a duas disjuntas. Removendo seus interiores, obtemos a variedade compacta com bordo
	\[
		X=M\setminus\bigcup_{i=1}^{\ell}\operatorname{int}(B_i).
	\]
	Escrevendo \(S_i=\partial B_i\), seu bordo é formado pelo bordo original de M e pelas esferas que estão nas redondezas das singularidades:
	\[
		\partial X=\partial M\sqcup\left(\bigsqcup_{i=1}^{\ell}(-S_i)\right),
	\]
	onde o sinal negativo indica a orientação oposta à orientação exterior usual de cada esfera \(S_i\).

	Como \(v\) não se anula em X, podemos considerar o campo normalizado
	\[
		\overline v:X\rightarrow\mathbb{S}^{m-1},\qquad
		\overline v(x)=\frac{v(x)}{\Vert v(x)\Vert};
	\]
	pelo lema sobre aplicações que se estendem a uma variedade com bordo, temos
	\[
		\mathrm{deg}(\overline v|_{\partial X})=0,
	\]
	e, levando em conta a orientação das componentes de \(\partial X\), obtemos
	\[
		0=\mathrm{deg}(\overline v|_{\partial M})
		-\sum_{i=1}^{\ell}\mathrm{deg}(\overline v|_{S_i}).
	\]
	Pela definição do índice de uma singularidade isolada,
	\[
		\mathrm{deg}(\overline v|_{S_i})=i(v,p_i),
	\]
	de onde segue que
	\[
		\boxed{\mathrm{deg}(\overline v|_{\partial M})
			=\sum_{i=1}^{\ell}i(v,p_i).}
	\]

	Por hipótese, o campo \(v\) aponta para fora ao longo do bordo, o que significa que, para \(x\in\partial M\), o ângulo entre \(v(x)\) e a normal exterior \(g(x)\) é estritamente menor que \(\frac{\pi}{2}\) (em particular, \(\overline v(x)\) e \(g(x)\) não são opostos).

	Podemos, então, construir uma homotopia entre \(g\) e \(\overline v|_{\partial M}\), normalizando a combinação dos dois campos:
	\[
		H(x,t)=\frac{g(x)\cos t+\overline v(x)\sin t}
		{\Vert g(x)\cos t+\overline v(x)\sin t\Vert},
		\qquad 0\leq t\leq\frac{\pi}{2}.
	\]
	Como o ângulo entre os dois vetores é agudo, o denominador é não nulo durante toda a homotopia. Além disso,
	\[
		H(x,0)=g(x),\qquad H\left(x,\frac{\pi}{2}\right)=\overline v(x).
	\]
	Pela invariância do grau por homotopia,
	\[
		\mathrm{deg}(\overline v|_{\partial M})=\mathrm{deg}(g).
	\]
	Portanto,
	\[
		\boxed{\sum_{j=1}^{\ell}i(v,p_j)=\mathrm{deg}(g).}
	\]
	(Continua...)
\end{proof*}
\end{document}

